\section*{Abstract}
\addcontentsline{toc}{section}{Abstract}
Even though artificial intelligence has achieved remarkable results over the past decade, several fundamental challenges remain unresolved. Current systems are highly energy intensive, struggle to learn continuously without interfering with existing knowledge, often fail to generalize reliably to novel situations, and lack the grounded causal understanding that comes from interacting with the world. If we want to address these problems, a natural place to look for inspiration is biology, and in particular the human brain, where many of these capabilities already exist. If future AI systems take stronger inspiration from biological architectures, they will also need learning rules that can operate within such architectures. Backpropagation trains today’s neural networks extremely well, but it is difficult to reconcile with biological learning. This motivates the search for biologically plausible learning algorithms.

This thesis studies perturbation-based learning, a branch of biologically plausible learning where random perturbations are added to a network’s weights or neural activity, and the resulting change in performance is used to decide whether these perturbations should be reinforced or reversed. The thesis first reviews the existing literature and identifies an unresolved gap in the relationship between the two main forms of perturbation learning. Building on an idea from Bayesian neural networks, it then addresses this gap theoretically by deriving a new result that clarifies this relationship. During this derivation, a new method called input-scaled node perturbation is introduced. To the best of our knowledge, neither the theoretical result nor the method has been derived before.

The empirical results provide evidence that these theoretical contributions also matter in practice. As the tested tasks become more complex, the relationship between the methods behaves in line with the theoretical prediction, and input-scaled node perturbation performs better than the existing perturbation methods it is compared against. These results should still be interpreted with caution, since the evaluation is limited to a specific set of supervised learning problems and network architectures. While the proposed method is promising relative to other perturbation methods, backpropagation remains far superior on the most challenging tasks studied here. The thesis therefore does not present perturbation-based learning as a replacement for backpropagation, but as a biologically motivated learning principle with real capability and room for further development. In particular, the results suggest that the theoretical direction developed in this thesis is a promising path for such further development.

\clearpage
